% !TEX root = ../paper.tex
% STATUS: No experiment below has been run. Protocols are written in final-paper tense.
% Every red \pending{} marker is a result, number, or claim heading that must be replaced by measured evidence or removed.
\section{Experimental Results}
We evaluate the proposed framework in both simulation and real-world experiments.
We first evaluate the simulation cost of moving-patch MPM against full-domain MPM and its ability to restore terrain deformation on revisits.
We then examine whether MPM facilitates successful policies across multiple embodiments and tasks.
Next, we analyze the contribution of MPM to locomotion policy performance by comparing policies trained on rigid, reduced-order, and MPM terrain.
We then deploy the humanoid policies on physical granular terrain.
Finally, we evaluate whether policies trained with Particle NeRD transfer to MPM and to hardware.

\subsection{Simulation Evaluation}
\noindent\textbf{Simulation Setup:}
All simulation experiments run on \pending{GPU}.
We use success rate (SR), velocity tracking error $e_v$, and yaw tracking error $e_\theta$ as metrics for humanoid locomotion.
The success rate is the fraction of \pending{100} episodes that complete the commanded course.
An episode fails if the base height falls below \pending{0.2}~m, if the roll or pitch magnitude exceeds \pending{$\pi/4$}~rad, or if the robot does not complete the course within \pending{10}~s.
The tracking errors are the mean absolute errors between commanded and measured base velocity and yaw rate.
Timing measurements include host-side tile paging and synchronization and exclude rendering and warm-up.

\noindent\textbf{\pending{Moving patches bound simulation cost independently of traversed distance.}}
We compare (1) full-domain MPM, (2) the recycling patch, and (3) the cached patch under matched material parameters, particle spacing, timestep, coupling, and robot trajectories.
We vary the patch size, number of parallel environments, and traversed distance, and measure peak GPU memory and coupled physics throughput.
We further compare the material reaction forces on the feet and the particle displacements near the robot against full-domain MPM, which serves as the numerical reference.
\pending{Result and interpretation.}

\noindent\textbf{\pending{Particle caching restores terrain deformation on revisits.}}
The robot walks a closed trajectory that returns to its own footprints.
We compare the terrain height and the reaction forces on the feet during the revisit among the recycling patch, the cached patch, and full-domain MPM.
\pending{Result and interpretation.}

\noindent\textbf{\pending{MPM facilitates successful policies across multiple embodiments and tasks.}}
We train policies for all three tasks of Sec.~\ref{sec:env_design}.
We evaluate Franka particle pushing by the fraction of material delivered to the target region and the spilled material, humanoid velocity tracking by SR, $e_v$, and $e_\theta$, and box pushing additionally by box velocity tracking and box tipping.
\pending{Result and interpretation.}
% TODO(figure): snapshots of the three tasks.

\noindent\textbf{\pending{MPM training improves locomotion on granular terrain.}}
To isolate the contribution of MPM, we compare locomotion policies trained on (1) rigid rough terrain in MuJoCo Warp (Rigid), (2) the 3D RFT contact model~\cite{gmloco} (RFT), (3) flat MPM terrain (MPM-Flat), and (4) rough MPM terrain (MPM-Rough).
All policies share the robot asset, observations, rewards, and training budget, and are evaluated on MPM terrain with held-out material parameters.
In addition to SR and tracking errors, we compare foot sinkage, gait timing, and pelvis acceleration to characterize the emergent gait.
Since evaluation in MPM favors MPM-trained policies, the hardware experiments in Sec.~\ref{sec:hardware} provide an independent comparison.
\pending{Result and interpretation.}

\subsection{Hardware Experiments}
\label{sec:hardware}
We deploy the trained policies on a Unitree G1 humanoid robot.
The policy outputs joint position targets at 50~Hz, which are sent to the low-level PD controller of the robot.
We test the policies on \pending{granular terrain types and bed depth}.
We consider an episode a failure if the robot fails to maintain an upright posture and requires human intervention.

\noindent\textbf{\pending{Claim for humanoid walking on physical granular terrain.}}
We compare Rigid, RFT, and MPM-Rough policies over commanded forward velocities of \pending{range}, with \pending{N} trials per command in randomized order.
\pending{Result and interpretation.}

\noindent\textbf{\pending{Claim for humanoid box pushing on physical granular terrain.}}
We track the box pose with motion capture and report SR, base velocity tracking, and box velocity tracking over \pending{N} trials.
\pending{Result and interpretation.}

\subsection{Particle NeRD Evaluation}
\noindent\textbf{\pending{Policies trained with Particle NeRD transfer to MPM.}}
We train Particle NeRD on a single terrain type and train a locomotion policy within the same material-parameter range.
We evaluate this policy in MPM with the same terrain type and material range, and compare it against a policy trained directly in MPM.
We report training wall time to reach \pending{threshold} together with the cost of data generation and network training.
Timing covers feature construction, network inference, contact reconstruction, rigid-body integration, and patch updates.
\pending{Result and interpretation.}

\noindent\textbf{\pending{Policies trained with Particle NeRD transfer to hardware.}}
\pending{Result and interpretation.}
